We now return to the question of existence. We will use what is called a \emph{direct method} in the calculus of variations. (We will try to minimize the energy directly rather than solve, for instance, the Euler-Lagrange equations.) The following lemma provides a first important step.

\begin{Lemma}
\label{lem:elasticity-weakly-convergent-minimizing-sequence}
There exists a minimizing sequence $\{\bm{p}_n\}\in\mathcal{D}_{\mathcal{E}}$ such that
\begin{equation}\label{eq:elasticity-minimizing-sequence-energy-infimum}
\lim_{n\to\infty}\mathcal{E}(\bm{p}_n)\leq\mathcal{E}(\bm{p})\quad\text{for all }\bm{p}\in\mathcal{D}_{\mathcal{E}}
\end{equation}
and with the additional property that $\bm{p}_n$ is weakly convergent; i.e., there exists $\widetilde{\bm{p}}\in W^{1,p}(\Omega)$ such that
\begin{equation}\label{eq:elasticity-minimizing-sequence-weak-limit}
\bm{p}_n\rightharpoonup\widetilde{\bm{p}}\text{ in }W^{1,p}(\Omega).
\end{equation}

\end{Lemma}

\begingroup
\renewcommand{\endProof}{\par\nobreak\hfill$\blacksquare$\par\vspace{0.1cm}}
\begin{Proof}
We first note that $\mathcal{E}$ is bounded below on $\mathcal{D}_{\mathcal{E}}$; i.e., for any $\bm{p}\in\mathcal{D}_{\mathcal{E}}$ we can use \eqref{eq:elasticity-stored-energy-p-growth-bound} and \eqref{eq:elasticity-potential-energy-uniform-lower-bound}
\begin{equation}\label{eq:elasticity-total-energy-lower-bound}
\mathcal{E}(\bm{p})\geq\int_\Omega\omega(\bm{x})+\Psi(\bm{p}(\bm{x}))\,\mathrm{d}\bm{x}\geq C-\|\omega\|_\infty|\Omega|,
\end{equation}
where $|\Omega|$ is the volume of $\Omega$. Since $\mathcal{E}(\bm{p})$ is bounded below it must have a greatest lower bound $L$, and hence there must be a sequence $\bm{p}_n\in\mathcal{D}_{\mathcal{E}}$ such that
\begin{equation}\label{eq:elasticity-minimizing-sequence-infimum-limit}
\lim_{n\to\infty}\mathcal{E}(\bm{p}_n)=L.
\end{equation}
Since $\mathcal{E}(\bm{p}_n)$ is a convergent sequence in $\mathbb{R}$ it must be bounded, say, by a constant $K$. Using this and \eqref{eq:elasticity-stored-energy-p-growth-bound} we get
\begin{equation}\label{eq:elasticity-minimizing-gradient-p-bound}
K\geq|\mathcal{E}(\bm{p}_n)|\geq k\int_\Omega|\nabla\bm{p}_n(\bm{x})|^p\,\mathrm{d}\bm{x}-\|\omega\|_\infty|\Omega|-C.
\end{equation}
Rearranging and combining this with Poincaré\textquotesingle s inequality (on $\bm{p}_n-\bm{g}$) gives us
\begin{equation}\label{eq:elasticity-minimizing-sobolev-norm-bound}
\|\bm{p}_n\|_{1,p}\leq\widetilde K,
\end{equation}
for some constant $\widetilde K$ independent of $n$.

Since $\bm{p}_n$ is bounded, Theorem~\ref{thm:alaoglu-weak-star-subsequence} implies that it has a weakly convergent subsequence (which we also label $\bm{p}_n$). Defining $\widetilde{\bm{p}}$ to be the weak limit of this sequence gives us \eqref{eq:elasticity-minimizing-sequence-weak-limit}. Since the original sequence $\mathcal{E}(\bm{p}_n)$ converges to $L$, so does any subsequence; this gives us \eqref{eq:elasticity-minimizing-sequence-energy-infimum} and completes the proof.
\end{Proof}
\endgroup

Of course, since $\bm{p}_n$ is a minimizing sequence, our first guess is that its ``limit,\textquotesingle\textquotesingle\ $\widetilde{\bm{p}}$, is a solution of our problem. However, two questions remain.\begin{enumerate}\item Is $\widetilde{\bm{p}}\in\mathcal{D}_{\mathcal{E}}$; i.e., is the constraint $\det\nabla\widetilde{\bm{p}}>0$ satisfied?\item If so, is $\widetilde{\bm{p}}$ actually a minimizer; i.e., is it true that
\begin{equation}\label{eq:elasticity-weak-limit-minimizer-question}
\mathcal{E}(\widetilde{\bm{p}})=\lim_{n\to\infty}\mathcal{E}(\bm{p}_n)=L\leq\mathcal{E}(\bm{p})\quad\text{for all }\bm{p}\in\mathcal{D}_{\mathcal{E}}?
\end{equation}
\end{enumerate}

In order to answer the first question, we extend the domain of definition of $\mathcal{E}$ to functions which do not satisfy the constraint \eqref{eq:elasticity-positive-orientation-determinant}. We define the function
\begin{equation}\label{eq:elasticity-extended-stored-density-mapping}
\Omega\times\mathcal{M}^{3\times3}\ni(\bm{x},\bm{f})\mapsto\overline{\mathcal{W}}(\bm{x},\bm{F})\in\overline{\mathbb{R}}:=\mathbb{R}\cup\{-\infty,\infty\}
\end{equation}
by
\begin{equation}\label{eq:elasticity-extended-density-by-determinant}
\overline{\mathcal{W}}(\bm{x},\bm{F}):=\begin{cases}\mathcal{W}(\bm{x},\bm{F}),&\det\bm{F}>0,\\\infty,&\det\bm{F}\leq0.\end{cases}
\end{equation}
We can now extend the domain of definition of our total energy to
\begin{equation}\label{eq:elasticity-extended-sobolev-affine-domain}
W_{\bm{g}}^{1,p}(\Omega):=\{\bm{p}\in W^{1,p}(\Omega)\mid\bm{p}-\bm{g}\in W_0^{1,p}(\Omega)\},
\end{equation}
by defining
\begin{equation}\label{eq:elasticity-extended-total-energy-functional}
\overline{\mathcal{E}}(\bm{p}):=\int_\Omega\overline{\mathcal{W}}(\bm{x},\bm{F}(\bm{x}))\,\mathrm{d}\bm{x}+\int_\Omega\Psi(\bm{p}(\bm{x}))\,\mathrm{d}\bm{x}.
\end{equation}
The following result is then immediate.

\begin{Lemma}
\label{lem:elasticity-finite-extended-energy-admissibility}
If $\overline{\mathcal{E}}(\bm{p})<\infty$, then $\bm{p}\in\mathcal{D}_{\mathcal{E}}$.
\end{Lemma}

This means that if we can indeed answer the second question in the affirmative; i.e., if $\widetilde{\bm{p}}$ is a minimizer, then $\widetilde{\bm{p}}\in\mathcal{D}_{\mathcal{E}}$. Thus, we focus on the second question and identify conditions that will ensure the weak limit is indeed a minimizer.

\begin{Definition}
\label{def:sequential-weak-semicontinuity-and-continuity}
Let $X$ be a Banach space. We say that a nonlinear mapping $\mathcal{F}:X\to\mathbb{R}$ is \textbf{sequentially weakly lower semicontinuous (wlsc)} if whenever
\begin{equation}\label{eq:sequential-weak-lower-semicontinuity-sequence}
v_n\rightharpoonup\overline v\text{ in }X,
\end{equation}
we have
\begin{equation}\label{eq:sequential-weak-lower-semicontinuity-inequality}
\mathcal{F}(\overline v)\leq\liminf_{n\to\infty}\mathcal{F}(v_n).
\end{equation}
We say that $\mathcal{F}:X\to\mathbb{R}$ is \textbf{sequentially weak-star lower semicontinuous (wslsc)} if whenever
\begin{equation}\label{eq:sequential-weak-star-lower-semicontinuity-sequence}
v_n\overset{*}{\rightharpoonup}\overline v\text{ in }X,
\end{equation}
it follows that \eqref{eq:sequential-weak-lower-semicontinuity-inequality} holds.

A mapping $\mathcal{F}$ is \textbf{sequentially weakly (weak-star) continuous} if
\begin{equation}\label{eq:sequential-weak-continuity-functional-limit}
\mathcal{F}(\overline v)=\lim_{n\to\infty}\mathcal{F}(v_n)
\end{equation}
whenever $v_n\rightharpoonup\overline v$ ($v_n\overset{*}{\rightharpoonup}\overline v$) in $X$.
\end{Definition}

\begin{Remark}
\label{rem:sequential-semicontinuity-terminology-omission}
We will drop the use of the word ``sequentially\textquotesingle\textquotesingle\ from here on. (A more thorough study of functional analysis would highlight the differences between the sequential and topological notions of continuity and lower semicontinuity.)
\end{Remark}

\begin{Remark}
\label{rem:fatou-lone-norm-lower-semicontinuity}
The reader already knows a basic lower semicontinuity result: \emph{Fatou\textquotesingle s lemma}, which implies (among other things) that the $L^1$ norm is lower semicontinuous.
\end{Remark}

We state a theorem (usually attributed to Tonelli) which is the fundamental result on weak lower semicontinuity. However, before doing this we give a definition of a convex function whose domain is in a general Banach space $X$ and whose range is the extended real line $\overline{\mathbb{R}}$.

\begin{Definition}
\label{def:extended-real-convex-function}
Let $K\subset X$ be a convex set. Then we say a mapping $G:K\to\overline{\mathbb{R}}$ is \textbf{convex} if for every $u,v\in K$, we have
\begin{equation}\label{eq:extended-real-convex-function-definition}
G(\lambda u+(1-\lambda)v)\leq\lambda G(u)+(1-\lambda)G(v)\quad\text{for all }\lambda\in[0,1],
\end{equation}
whenever the right-hand side of the inequality is well defined.
\end{Definition}

\begin{Theorem}{Tonelli}
\label{thm:tonelli-convexity-weak-lower-semicontinuity}
For functions $\bm{u}:\Omega\to\mathbb{R}^m$, define the nonlinear function
\begin{equation}\label{eq:tonelli-integral-functional-definition}
\mathcal{F}(\bm{u}):=\int_\Omega f(\bm{u}(\bm{x}))\,\mathrm{d}\bm{x},
\end{equation}
where $f:\mathbb{R}^m\to\overline{\mathbb{R}}$ is continuous. Then the function $\mathcal{F}$ is sequentially weakly lower semicontinuous on the space $L^p(\Omega)$ for $1<p<\infty$ and weak-star lower semicontinuous on $L^\infty(\Omega)$ if and only if $\mathbb{R}^m\ni\bm{u}\mapsto f(\bm{u})\in\overline{\mathbb{R}}$ is convex.
\end{Theorem}

\begingroup
\renewcommand{\endProof}{\par}
\begin{Proof}
The proof of this theorem is found in many texts on convex analysis or the calculus of variations. We give only a sketch of the proof here. We begin by showing that weak-star lower semicontinuity of $\mathcal{F}$ in $L^\infty(\Omega)$ implies that $f$ is convex. We first prove a lemma, which highlights one of the most important types of weak convergence: a wildly oscillating sequence which converges to its average value.
\end{Proof}
\endgroup

\begin{Lemma}
\label{lem:oscillating-sequence-average}
For any $\bm{a},\bm{b}\in\mathbb{R}^m$ and $\theta\in(0,1)$ define $\bm{u}:\mathbb{R}\to\mathbb{R}^m$ to be a function of period one such that
\begin{equation}\label{eq:oscillating-two-value-periodic-function}
\bm{u}(x):=\begin{cases}\bm{a},&x\in[0,\theta),\\\bm{b},&x\in[\theta,1).\end{cases}
\end{equation}
Let the sequence of functions $\bm{u}_n:[0,1]\to\mathbb{R}^m$ be given by
\begin{equation}\label{eq:oscillating-rescaled-periodic-sequence}
\bm{u}_n(x):=\bm{u}(nx),\quad x\in[0,1].
\end{equation}
Then
\begin{equation}\label{eq:oscillating-sequence-weak-lp-limit}
\bm{u}_n\rightharpoonup\overline{\bm{u}}
\end{equation}
in $L^p(0,1)$ for $1<p<\infty$ and
\begin{equation}\label{eq:oscillating-sequence-weak-star-linfty-limit}
\bm{u}_n\overset{*}{\rightharpoonup}\overline{\bm{u}}
\end{equation}
in $L^\infty(0,1)$, where $\overline{\bm{u}}:=\theta\bm{a}+(1-\theta)\bm{b}$ is a constant function.
\end{Lemma}

\begingroup
\renewcommand{\endProof}{\par\nobreak\hfill$\blacksquare$\par\vspace{0.1cm}}
\begin{Proof}
Here as well, we just give a sketch of the proof. We wish to show that
\begin{equation}\label{eq:oscillating-sequence-tested-average-limit}
\lim_{n\to\infty}\int_0^1\bm{u}_n(x)\phi(x)\,\mathrm{d}x=[\theta\bm{a}+(1-\theta)\bm{b}]\int_0^1\phi(x)\,\mathrm{d}x,
\end{equation}
for any $\phi\in L^p(0,1)$, $1\leq p<\infty$. The reader should fill in the following steps.\begin{enumerate}\item The assertion is easy to prove for functions of the form
\begin{equation}\label{eq:oscillating-sequence-interval-step-test}
\phi(x)=\begin{cases}c,&x\in I\\0,&x\notin I,\end{cases}
\end{equation}
where $I\subset[0,1]$ is an interval. (For large $n$ the interval $I$ will contain a large integral number of periods of $\bm{u}_n$ plus some small ``slop\textquotesingle\textquotesingle\ at the ends.)\item One can use the previous observation to show that \eqref{eq:oscillating-sequence-tested-average-limit} holds if $\phi$ is a \emph{simple function}; i.e., if $\phi$ is piecewise constant with a finite number of jump discontinuities.\item One can then show that for any $\phi\in L^p(0,1)$, $p\in[1,\infty)$, and any $\epsilon>0$, we can find a simple function $\phi_s$ such that
\begin{equation}\label{eq:oscillating-sequence-simple-function-lone-approximation}
\|\phi-\phi_s\|_1<\epsilon.
\end{equation}
(By the definition of $L^1$ given in this book, we can approximate $\phi$ arbitrarily closely by a bounded, continuous function. The reader should verify that a bounded continuous function can be approximated arbitrarily closely in the $L^1$ norm by a simple function.)\end{enumerate}The combination of these observations completes the proof.
\end{Proof}
\endgroup

\begin{Remark}
\label{rem:oscillating-sequence-interval-and-dimension}
The choice of the interval $[0,1]$ was arbitrary. With only trivial modifications of the proof, one can show that
\begin{equation}\label{eq:oscillating-sequence-general-compact-interval-limit}
\bm{u}_n\rightharpoonup\overline{\bm{u}}
\end{equation}
in $L^p(I)$, for any compact interval $I\subset\mathbb{R}$. Furthermore, we proved this result for functions whose domains lie in a subset of $\mathbb{R}$ only for clarity; an analogous construction can be created using functions $\bm{u}_n$ defined on the domain $\Omega\subset\mathbb{R}^n$ by letting the functions oscillate in a single coordinate direction.
\end{Remark}

We now return to the proof of Tonelli\textquotesingle s theorem. We assume that $\mathcal{F}$ is weak-star lower semicontinuous. Let $C\subset\Omega\subset\mathbb{R}^n$ be a hypercube. For any $\bm{a},\bm{b}\in\mathbb{R}^m$, $\theta\in(0,1)$, let $\bm{u}_n$ be a sequence of oscillating functions with support on $C$ such that
\begin{equation}\label{eq:tonelli-necessity-two-state-weak-limit}
\bm{u}_n\rightharpoonup\overline{\bm{u}}=\theta\bm{a}+(1-\theta)\bm{b}
\end{equation}
in $L^\infty(\Omega)$. The key observation here is that the sequence of composite functions $f(\bm{u}_n)$ oscillate between the values $f(\bm{a})$ and $f(\bm{b})$ with volume fractions $\theta$ and $1-\theta$, respectively. Hence, it follows from the arguments in the proof of the lemma above that
\begin{equation}\label{eq:tonelli-necessity-composite-weak-star-limit}
f(\bm{u}_n)\overset{*}{\rightharpoonup}\theta f(\bm{a})+(1-\theta)f(\bm{b})
\end{equation}
in $L^\infty(\Omega)$. Combining this with the weak-star lower semicontinuity of $\mathcal{F}$, we have\[
\begin{aligned}
 |\Omega|f(\theta\bm{a}+(1-\theta)\bm{b})&=\mathcal{F}(\overline{\bm{u}})\\
 &\leq\liminf_{n\to\infty}\mathcal{F}(\bm{u}_n)\\
 &=\lim_{n\to\infty}\int_\Omega f(\bm{u}_n)\\
 &=|\Omega|\{\theta f(\bm{a})+(1-\theta)f(\bm{b})\}.
\end{aligned}
\]It follows that $f$ is convex.

We now assume that $f$ is convex and show that $\mathcal{F}$ is weak-star lower semicontinuous. We let $\bm{u}_n\rightharpoonup\overline{\bm{u}}$ be an arbitrary weakly convergent sequence and let
\begin{equation}\label{eq:tonelli-sufficiency-liminf-energy}
L=\liminf\mathcal{F}(\bm{u}_n).
\end{equation}
By taking a subsequence, we can actually assume that
\begin{equation}\label{eq:tonelli-sufficiency-subsequence-energy-limit}
L=\lim\mathcal{F}(\bm{u}_n).
\end{equation}


To complete the proof of the theorem we will use Mazur\textquotesingle s lemma which we state without proof. (The proof uses elementary convex analysis and is found in many texts (cf., e.g., \cite{ET}).)

\begin{Lemma}{Mazur}
\label{lem:mazur-strong-convex-combinations}
Let $X$ be a Banach space and suppose
\begin{equation}\label{eq:mazur-weakly-convergent-sequence}
u_n\rightharpoonup\overline u
\end{equation}
in $X$. Then there exists a function $N:\mathbb{N}\to\mathbb{N}$, and a sequence of sets of real numbers $\{\alpha(n)_k\}_{k=n}^{N(n)}$ such that $\alpha(n)_k\geq0$ and $\sum_{k=n}^{N(n)}\alpha(n)_k=1$ such that the sequence
\begin{equation}\label{eq:mazur-strongly-convergent-convex-combinations}
v_n:=\sum_{k=n}^{N(n)}\alpha(n)_ku_k
\end{equation}
converges strongly to $\overline u$ in $X$.
\end{Lemma}

\begin{Remark}
\label{rem:convex-combination-and-convex-hull}
We say that $v_n$ as defined above is a \emph{convex combination} of elements of the set $\{u_k\}_{k=n}^{N(n)}$. The set of all possible convex combinations of elements of a set $S$ is called the \emph{convex hull} of $S$.
\end{Remark}

\begingroup
\renewenvironment{Proof}{\par}{\par\nobreak\hfill$\blacksquare$\par\vspace{0.1cm}}
\begin{Proof}
We can use Mazur\textquotesingle s lemma to construct a sequence such that $\bm{v}_n\to\overline{\bm{u}}$ strongly in $L^p(\Omega)$ for every $p\in[1,\infty)$. Thus, at least for a subsequence we have $\bm{v}_n\to\overline{\bm{u}}$ almost everywhere. There are two steps to the remainder of the proof of Theorem~\ref{thm:tonelli-convexity-weak-lower-semicontinuity}.\begin{enumerate}\item Since $f(\bm{v}_n)\to f(\overline{\bm{u}})$ almost everywhere, it follows from Fatou\textquotesingle s lemma and the convexity of $f$ that
\begin{equation}\label{eq:tonelli-sufficiency-fatou-lower-bound}
\mathcal{F}(\overline{\bm{u}})\leq\liminf_{n\to\infty}\mathcal{F}(\bm{v}_n).
\end{equation}
(If $f$ is bounded below, as it is in most applications, we can use Fatou\textquotesingle s lemma without using convexity. Otherwise, one can use the fact that any convex function is bounded below by an affine function.)\item We now let $\epsilon>0$ be given and choose $N$ sufficiently large so that $\mathcal{F}(\overline{\bm{u}})\leq\mathcal{F}(\bm{v}_k)+\epsilon/2$ and $\mathcal{F}(\bm{u}_k)\leq L+\epsilon/2$ for $k\geq N$. Then for $n\geq N$ we have\[
\begin{aligned}
 \mathcal{F}(\overline{\bm{u}})&\leq\mathcal{F}(\bm{v}_n)+\epsilon/2\\
 &=\int_\Omega f\left(\sum_{k=n}^{N(n)}\alpha(n)_k\bm{u}_k\right)\,\mathrm{d}\bm{x}+\epsilon/2\\
 &\leq\sum_{k=n}^{N(n)}\alpha(n)_k\int_\Omega f(\bm{u}_k)\,\mathrm{d}\bm{x}+\epsilon/2\\
 &\leq L+\epsilon.
\end{aligned}
\]Here we have used the fact that $\alpha(n)_k\geq0$ and $\sum_{k=n}^{N(n)}\alpha(n)_k=1$. Since $\epsilon$ was arbitrary, the proof is complete.\end{enumerate}
\end{Proof}
\endgroup

By applying the previous theorem to $\mathcal{F}$ and $-\mathcal{F}$ we immediately get the following result.

\begin{Corollary}
\label{cor:tonelli-affine-integrand-weak-continuity}
The mapping $\mathcal{F}$ is weakly (weak-star) continuous if and only if $\bm{u}\mapsto f(\bm{u})$ is affine; i.e.,
\begin{equation}\label{eq:weak-continuity-integral-affine-integrand}
f(\bm{u})=\alpha+\bm{b}\cdot\bm{u},
\end{equation}
for some $\alpha\in\mathbb{R}$ and some $\bm{b}\in\mathbb{R}^n$.
\end{Corollary}

(The reader should verify (or already know) that if both $f$ and $-f$ are convex, then $f$ is affine.)

Tonelli\textquotesingle s theorem can now be applied directly to energy functional $\overline{\mathcal{E}}$.

\begin{Corollary}
\label{cor:convex-elasticity-energy-equilibrium-existence}
If the function $\bm{F}\mapsto\overline{\mathcal{W}}(\bm{x},\bm{F})$ is convex, then $\overline{\mathcal{E}}$ is weakly lower semicontinuous. Furthermore, there exists an equilibrium state $\overline{\bm{p}}$.
\end{Corollary}

\begingroup
\renewcommand{\endProof}{\par\nobreak\hfill$\blacksquare$\par\vspace{0.1cm}}
\begin{Proof}
The theorem applies directly to the minimizing sequence and the stored energy $\mathcal{W}$. To take care of the potential energy term we use compact imbedding to get $\bm{p}_n\to\overline{\bm{p}}$ strongly in $L^p(\Omega)$. Then, using Fatou\textquotesingle s lemma once again, we get
\begin{equation}\label{eq:elasticity-potential-energy-fatou-weak-limit}
\int_\Omega\Psi(\overline{\bm{p}})\,\mathrm{d}\bm{x}\leq\liminf_{n\to\infty}\int_\Omega\Psi(\bm{p}_n)\,\mathrm{d}\bm{x}
\end{equation}
This completes the proof.
\end{Proof}
\endgroup

The result above is useful in some situations (e.g., linear elasticity). However, in problems in nonlinear elasticity (where we really want to apply this theory in the first place) there are good reasons why the assumption of convexity of the energy density is physically unreasonable. We will not discuss these reasons here; for an introductory discussion of these issues the reader could consult \cite{Ba}.

If we weaken the convexity assumption on $\mathcal{W}$, Tonelli\textquotesingle s theorem would seem to imply that we will not be able to show that the energy $\mathcal{E}$ is weakly lower semicontinuous. However, this is not the case. We know more about sequences $\bm{F}_n\rightharpoonup\overline{\bm{F}}$ than the fact that they are weakly convergent. We know that each element of the sequence is a gradient; i.e., $\bm{F}_n=\nabla\bm{p}_n$. How can we identify which matrix-valued functions $\bm{F}:\mathbb{R}^3\to Q$ are given by gradients? Using the equality of mixed partial derivatives, we see that
\begin{equation}\label{eq:deformation-gradient-mixed-partial-constraint}
\dfrac{\partial F_{ij}}{\partial x_k}=\dfrac{\partial^2p_i}{\partial x_j\partial x_k}=\dfrac{\partial^2p_i}{\partial x_k\partial x_j}=\dfrac{\partial F_{ik}}{\partial x_j}.
\end{equation}
(This is clear if $\bm{p}$ is smooth. It holds in the space $W^{-1,p}(\Omega)$ for general $\bm{p}\in W^{1,p}(\Omega)$.) This is simply a complicated version of the condition that the curl of a gradient be zero. With this in mind, it is natural to study weak lower semicontinuity under the assumption of \emph{differential constraints}. A very important example of a result of this type is the following: often called the \emph{div-curl lemma}.

\begin{Theorem}{Div-Curl}
\label{thm:div-curl-compact-negative-norm}
Let $\Omega\subseteq\mathbb{R}^n$ be a domain. Suppose $\bm{u}^k:\Omega\to\mathbb{R}^n$ and $\bm{v}^k:\Omega\to\mathbb{R}^n$ are sequences of vector valued functions satisfying
\begin{equation}\label{eq:div-curl-first-vector-weak-limit}
\bm{u}^k\rightharpoonup\overline{\bm{u}}\text{ in }L^2(\Omega),
\end{equation}

\begin{equation}\label{eq:div-curl-second-vector-weak-limit}
\bm{v}^k\rightharpoonup\overline{\bm{v}}\text{ in }L^2(\Omega),
\end{equation}
and suppose the sequences
\begin{equation}\label{eq:div-curl-divergence-sequence}
\operatorname{div}\bm{u}^k:=\sum_{i=1}^n\dfrac{\partial u_i^k}{\partial x_i},
\end{equation}

\begin{equation}\label{eq:div-curl-curl-sequence}
\operatorname{curl}\bm{v}^k:=\dfrac{\partial v_j^k}{\partial x_i}-\dfrac{\partial v_i^k}{\partial x_j},\quad i,j=1,\ldots,n,
\end{equation}
lie in a compact set in $H_{\mathrm{loc}}^{-1}(\Omega)$. Then
\begin{equation}\label{eq:div-curl-product-distributional-limit}
\bm{u}^k\cdot\bm{v}^k\to\overline{\bm{u}}\cdot\overline{\bm{v}}\text{ in }\mathcal{D}^{\prime}(\Omega).
\end{equation}

\end{Theorem}

Here, $H_{\mathrm{loc}}^{-1}(\Omega)$ is the set of functions $f$ such that for any test function $\phi\in\mathcal{D}(\Omega)$ we have $\phi f\in H^{-1}(\Omega)$.

Since we won\textquotesingle t use this lemma below we will skip the proof. However, we leave the proof of an easier version as an exercise (cf. Problem~\ref{ex:div-curl-scalar-gradient-version}).

A more useful result from the point of view of our problems in elasticity is the following theorem on the weak continuity of subdeterminants of gradients.

\begin{Theorem}
\label{thm:gradient-minors-weak-continuity}
Let $\Omega\subset\mathbb{R}^n$ be a bounded domain. Suppose $n<p<\infty$ and suppose the sequence of functions $\bm{p}_k:\Omega\to\mathbb{R}^n$ satisfy
\begin{equation}\label{eq:gradient-minors-weak-sobolev-sequence}
\bm{p}_k\rightharpoonup\overline{\bm{p}}\text{ in }W^{1,p}(\Omega).
\end{equation}
Let $m\leq n$ and let $M_k$ be the sequence of $m\times m$ subdeterminants obtained by taking a fixed $m$ rows and $m$ columns of $\nabla\bm{p}_k$, and let $\overline M$ be the corresponding subdeterminant of $\nabla\overline{\bm{p}}$. Then
\begin{equation}\label{eq:gradient-minors-weak-lp-over-order-limit}
M_k\rightharpoonup\overline M\quad\text{in }L^{p/m}(\Omega).
\end{equation}

\end{Theorem}

\begingroup
\renewcommand{\endProof}{\par\nobreak\hfill$\blacksquare$\par\vspace{0.1cm}}
\begin{Proof}
The proof proceeds by induction, and the first step is left to the reader (Problem~\ref{ex:gradient-minors-induction-first-step}).

Without loss of generality, we can complete the proof by showing that
\begin{equation}\label{eq:gradient-determinant-weak-lp-over-dimension-limit}
\det(\nabla\bm{p}_k)\rightharpoonup\det(\nabla\overline{\bm{p}})\quad\text{in }L^{p/n}(\Omega)
\end{equation}
under the assumption that any $(n-1)\times(n-1)$ subdeterminant $M_k$ satisfies \eqref{eq:gradient-minors-weak-lp-over-order-limit} (with $m=n-1$). However, under this assumption we have
\begin{equation}\label{eq:gradient-cofactor-weak-limit-induction}
\operatorname{cof}\nabla\bm{p}_k\rightharpoonup\operatorname{cof}\nabla\overline{\bm{p}}\quad\text{in }L^{p/(n-1)}(\Omega).
\end{equation}
Here we have used $\operatorname{cof}\bm{A}$ to denote the \emph{cofactor matrix} of $\bm{A}$. (The $(i,j)$th component of the cofactor matrix is $(-1)^{i+j}$ times the $(i,j)$th minor.)

We now use the fact that for smooth functions $\bm{p}$
\begin{equation}\label{eq:gradient-determinant-cofactor-row-expansion}
\det\nabla\bm{p}=\sum_j\dfrac{\partial p_i}{\partial x_j}(\operatorname{cof}\nabla\bm{p})_{ij},
\end{equation}
for any $i=1,\ldots,n$. In addition, we use the identity
\begin{equation}\label{eq:gradient-cofactor-row-divergence-identity}
\sum_j\dfrac{\partial}{\partial x_j}(\operatorname{cof}\nabla\bm{p})_{ij}=0,\quad i=1,\ldots,n.
\end{equation}
(This identity can be verified by direct calculation (cf. Problem~\ref{ex:subdeterminants-null-lagrangians}).) Therefore, we can derive the formula
\begin{equation}\label{eq:gradient-determinant-divergence-form}
\det\nabla\bm{p}=\sum_j\dfrac{\partial}{\partial x_j}\bigl(p_i(\operatorname{cof}\nabla\bm{p})_{ij}\bigr).
\end{equation}
Thus, after approximating our sequence $\bm{p}_k$ by smooth functions, we have for any $\phi\in\mathcal{D}(\Omega)$\[
\begin{aligned}
 \int_\Omega\phi\det\nabla\bm{p}_k\,\mathrm{d}\bm{x}&=-\int_\Omega\sum_j\dfrac{\partial\phi}{\partial x_j}p_i^k(\operatorname{cof}\nabla\bm{p}_k)_{ij}\,\mathrm{d}\bm{x}\\
 &\to-\int_\Omega\sum_j\dfrac{\partial\phi}{\partial x_j}\overline p_i(\operatorname{cof}\nabla\overline{\bm{p}})_{ij}\,\mathrm{d}\bm{x}\\
 &=\int_\Omega\phi\det\nabla\overline{\bm{p}}\,\mathrm{d}\bm{x}.
\end{aligned}
\]In taking the limit above we have used the fact that by compact imbedding $\bm{p}^k\to\overline{\bm{p}}$ (strongly) in $L^p(\Omega)$, and hence (since $\Omega$ is bounded) $\bm{p}^k\to\overline{\bm{p}}$ (strongly) in $L^q(\Omega)$ where $q=p/(1+p-n)$ is the conjugate exponent of $p/(n-1)$. We can use this and Problem~\ref{ex:weak-strong-holder-dual-product-distribution-limit} to take the limit.

Thus, $\det\nabla\bm{p}_k\to\det\nabla\overline{\bm{p}}$ in $\mathcal{D}^{\prime}(\Omega)$. To complete the proof, we use the fact that $\det\nabla\bm{p}_k$ is bounded in $L^{p/n}(\Omega)$ and a density argument (cf. Problem~\ref{ex:bounded-l-two-distribution-limit-is-weak-limit}).
\end{Proof}
\endgroup

We can use this result coupled with the following definition in our study of variational problems of nonlinear elasticity.

\begin{Definition}
\label{def:polyconvex-function}
Let $\mathcal{M}^{m\times n}$ be the set of $m\times n$ matrices. A function $G:\mathcal{M}^{m\times n}\to\mathbb{R}$ is said to be \textbf{polyconvex} if $\bm{A}\mapsto G(\bm{A})$ can be represented as a convex function of the subdeterminants of $\bm{A}$.
\end{Definition}

In the particular case of three-dimensional elasticity, we say that the stored energy density $\mathcal{W}(\bm{x},\bm{F})$ is polyconvex if there exists a function
\begin{equation}\label{eq:polyconvex-elastic-energy-convex-representation-map}
\Omega\times\mathcal{M}^{3\times3}\times\mathcal{M}^{3\times3}\times(0,\infty)\ni(\bm{x},\bm{A},\bm{B},d)\mapsto g(\bm{x},\bm{A},\bm{B},d)\in\mathbb{R},
\end{equation}
such that for every $\bm{x}\in\Omega$, $(\bm{A},\bm{B},d)\mapsto g(\bm{x},\bm{A},\bm{B},d)$ is convex and
\begin{equation}\label{eq:polyconvex-elastic-energy-gradient-minor-representation}
\mathcal{W}(\bm{x},\bm{F})=g(\bm{x},\bm{F},\operatorname{cof}\bm{F},\det\bm{F}).
\end{equation}


\begin{Remark}
\label{rem:polyconvexity-weaker-than-convexity}
Polyconvexity is indeed a weaker assumption than convexity. In particular, such functions as
\begin{equation}\label{eq:polyconvex-nonconvex-inverse-determinant-example}
G(\bm{F})=\begin{cases}\dfrac{1}{\det\bm{F}},&\det\bm{F}>0,\\\infty,&\det\bm{F}\leq0,\end{cases}
\end{equation}
are polyconvex but not convex.
\end{Remark}

\begin{Remark}
\label{rem:deformation-minors-geometric-meaning}
In elasticity, where $\bm{F}(\bm{x})$ is the deformation gradient at $\bm{x}\in\Omega$, the components of $\bm{F}(\bm{x})$ reflect local changes under the deformation in the length of curves going through $\bm{x}$, the components of $\operatorname{cof}\bm{F}(\bm{x})$ reflect changes in the areas of variously oriented surfaces and $\det\bm{F}(\bm{x})$ reflects the local change in volume.
\end{Remark}

\begin{Remark}
\label{rem:polyconvexity-physical-reasonableness}
The physical objections raised in the literature for convex stored energy functions do not apply to polyconvex functions. Polyconvexity is widely accepted as an assumption which is both extremely general and physically reasonable.
\end{Remark}

We now state our basic existence result for elastic materials with polyconvex energies.

\begin{Theorem}
\label{thm:polyconvex-energy-equilibrium-existence}
Let $\overline{\mathcal{W}}$ satisfy assumptions 1--5 on page~\pageref{eq:elasticity-stored-energy-p-growth-bound} and, in addition, be polyconvex. Then there exists a minimizer $\widetilde{\bm{p}}\in\mathcal{D}_{\mathcal{E}}$ of the energy $\overline{\mathcal{E}}$.
\end{Theorem}

\begingroup
\renewcommand{\endProof}{\par\nobreak\hfill$\blacksquare$\par\vspace{0.1cm}}
\begin{Proof}
By Lemma~\ref{lem:elasticity-weakly-convergent-minimizing-sequence} there exists a minimizing sequence $\bm{p}_n\in\mathcal{D}_{\mathcal{E}}$ and a function $\widetilde{\bm{p}}\in W_{\bm{g}}^{1,p}(\Omega)$ such that $\bm{p}_n\rightharpoonup\widetilde{\bm{p}}$ in $W^{1,p}(\Omega)$. Now let $\bm{F}_n:=\nabla\bm{p}_n$, $\widetilde{\bm{F}}:=\nabla\widetilde{\bm{p}}$. Then Theorem~\ref{thm:gradient-minors-weak-continuity} gives us
\begin{equation}\label{eq:polyconvex-existence-cofactor-weak-limit}
\operatorname{cof}\bm{F}_n\rightharpoonup\operatorname{cof}\widetilde{\bm{F}}\quad\text{in }L^{p/2}(\Omega),
\end{equation}

\begin{equation}\label{eq:polyconvex-existence-determinant-weak-limit}
\det\bm{F}_n\rightharpoonup\det\widetilde{\bm{F}}\quad\text{in }L^{p/3}(\Omega).
\end{equation}
Now, we use Theorem~\ref{thm:tonelli-convexity-weak-lower-semicontinuity} to get\[
\begin{aligned}
\overline{\mathcal{E}}(\widetilde{\bm{p}})&=\int_\Omega\overline{\mathcal{W}}(\bm{x},\widetilde{\bm{F}}(\bm{x}))+\Psi(\widetilde{\bm{p}}(\bm{x}))\,\mathrm{d}\bm{x}\\
&=\int_\Omega g(\bm{x},\widetilde{\bm{F}}(\bm{x}),\operatorname{cof}\widetilde{\bm{F}}(\bm{x}),\det\widetilde{\bm{F}}(\bm{x}))+\Psi(\widetilde{\bm{p}}(\bm{x}))\,\mathrm{d}\bm{x}\\
&\leq\liminf_{n\to\infty}\int_\Omega g(\bm{x},\bm{F}_n(\bm{x}),\operatorname{cof}\bm{F}_n(\bm{x}),\det\bm{F}_n(\bm{x}))+\Psi(\bm{p}_n(\bm{x}))\,\mathrm{d}\bm{x}\\
&=\lim_{n\to\infty}\int_\Omega\overline{\mathcal{W}}(\bm{x},\bm{F}_n(\bm{x}))+\Psi(\bm{p}_n(\bm{x}))\,\mathrm{d}\bm{x}\\
&=L.
\end{aligned}
\]Where $L$ is the greatest lower bound of $\overline{\mathcal{E}}$ over $\mathcal{D}_{\mathcal{E}}$. Thus, $\widetilde{\bm{p}}$ is indeed a minimizer and this completes the proof.
\end{Proof}
\endgroup
